\section{Checkpoint, Plan, and Configuration Formats}
\label{app:formats}

The formats are designed for bounded range access, integrity, portability, and streaming conversion. JSON examples are illustrative; the repository includes normative JSON Schemas.

\subsection{AMS package layout}

A directory package has the form:

\begin{lstlisting}
model.ams/
  manifest.json
  graph.amsir
  tensors/
    shard-00000.amsbin
    shard-00001.amsbin
  packs/
    cpu-generic/...
    cuda-smXX/...
  signatures/
    manifest.ed25519
  provenance/
    source.json
\end{lstlisting}

A single-file container may place a fixed header, canonical manifest, chunk table, aligned data chunks, and optional signature footer in one file. The directory form is simpler for updates and content-addressed stores. All offsets are unsigned 64-bit with checked arithmetic; implementations may impose lower configured limits.

\subsection{Tensor descriptor}

Each tensor descriptor separates logical shape and alias semantics from one or more independently addressable physical layouts.

\begin{lstlisting}[language={},caption={Illustrative tensor manifest entry.}]
{
  "tensor_id": "model.layers.0.mlp.down_proj.weight",
  "storage_id": "sha256:...",
  "shape": [4096, 11008],
  "logical_dtype": "bfloat16",
  "storage_dtype": "int4_grouped",
  "byte_order": "little",
  "strides": [11008, 1],
  "storage_offset_elements": 0,
  "immutable": true,
  "aliases": [],
  "codec": {
    "name": "ams.int4.grouped.v1",
    "group_axis": 1,
    "group_size": 128,
    "scale_dtype": "float16",
    "zero_point": false
  },
  "layouts": [
    {
      "layout_id": "generic-2d-256x2048-v1",
      "tile_shape": [256, 2048],
      "chunks": [
        {
          "object": "tensors/shard-00000.amsbin",
          "offset": 1048576,
          "length": 278528,
          "logical_origin": [0, 0],
          "logical_extent": [256, 2048],
          "checksum": "sha256:..."
        }
      ]
    }
  ]
}
\end{lstlisting}

Every logical tile is covered exactly once per complete layout, except declared sparse/default regions. Chunks cannot overlap reserved metadata ranges. Codec metadata has a declared maximum decoded size and validation function.

\subsection{Root manifest}

The root records:

\begin{itemize}
  \item format major/minor version and feature flags;
  \item model architecture/configuration and canonical graph hash;
  \item tokenizer and generation asset hashes;
  \item tensor descriptors and unique storage objects;
  \item aliases, views, tied weights, and mutable state roots;
  \item operator/plugin requirements and minimum versions;
  \item source provenance, conversion command, and licenses;
  \item content hashes, optional Merkle root, and signature metadata;
  \item creation time as metadata only, never as part of deterministic content hashing unless declared.
\end{itemize}

Canonical JSON uses sorted keys, UTF-8, normalized numbers, and no insignificant whitespace for signing. Alternatively, a canonical binary format may be used; readers must not sign an implementation-dependent serialization.

\subsection{Plan format}

A plan document includes:

\begin{lstlisting}[language={},caption={Abbreviated plan structure.}]
{
  "plan_schema": "1.0",
  "plan_id": "sha256:...",
  "model_root": "sha256:...",
  "graph_hash": "sha256:...",
  "hardware_fingerprint": "sha256:...",
  "software_fingerprint": "sha256:...",
  "request_envelope": {
    "max_prompt_tokens": 4096,
    "max_new_tokens": 512,
    "max_batch": 4
  },
  "numeric_mode": "deterministic",
  "reservations": {
    "device:0": 6442450944,
    "host:pinned:0": 268435456,
    "host:cache:0": 8589934592,
    "spill": 274877906944
  },
  "operators": [
    {
      "node_id": 317,
      "implementation_id": "ams.linear.cuda.v3",
      "tile": {"p": 64, "i": 512, "j": 2048},
      "loop_order": "output_stationary",
      "buffers": 2,
      "workspace_max": 8388608,
      "resource_vector": {"device:0": 75497472}
    }
  ],
  "fallback_plan_ids": ["sha256:..."],
  "calibration_id": "sha256:..."
}
\end{lstlisting}

Plan IDs hash all semantics- and resource-relevant fields. Human-readable estimates may be excluded from the hash if clearly marked non-normative. Plans are caches: rejection after a software/hardware mismatch triggers regeneration, not model corruption.

\subsection{Runtime configuration schema}

The normative schema is closed by default. A minimal configuration contains:

\begin{lstlisting}[language={},caption={Illustrative runtime configuration.}]
{
  "version": "1.0",
  "storage": {
    "model_uri": "/models/model.ams",
    "spill_uri": "/var/lib/ams/spill",
    "spill_quota_bytes": 300000000000,
    "direct_io": "auto",
    "verify_checksums": true
  },
  "memory": {
    "device_budget_bytes": {"device:0": 6442450944},
    "pinned_budget_bytes": 268435456,
    "host_cache_budget_bytes": 8589934592,
    "runtime_reserve_fraction": 0.08
  },
  "execution": {
    "backend_preference": ["cuda", "cpu"],
    "numeric_mode": "stable",
    "max_inflight_requests": 4,
    "deterministic": false,
    "allow_approximation": false
  },
  "security": {
    "allow_pickle": false,
    "require_signature": false,
    "allowed_plugins": []
  },
  "telemetry": {
    "metrics": true,
    "trace_sampling": 0.01,
    "log_tensor_data": false
  }
}
\end{lstlisting}

The reserve fraction is a planning input, not an allocation guarantee. The broker converts it into a concrete reservation or returns preflight failure.

\subsection{Journals and version roots}

A journal record has transaction ID, monotonically increasing sequence, record type, logical tile/version, physical region, checksum, and checksum of the record. Publication writes a root record naming a complete set or parent-plus-delta chain. Recovery scans valid records, ignores torn tails, and selects the newest atomically published root whose referenced regions verify.

\subsection{Integrity and signatures}

Chunk checksums detect corruption; a signed root authenticates the mapping from logical tensors to chunk hashes. Signatures do not encrypt weights. Deployments may encrypt storage objects through filesystem, block-device, or store adapters; direct device IO must preserve authentication/decryption semantics and resource bounds.

\subsection{Alignment and portability}

Generic chunks use a conservative alignment, such as 4~KiB, and contain explicit padding excluded from logical data. Device-specific packs state target architecture, compiler/kernel ABI, alignment, and source logical hash. A reader may ignore incompatible packs and use the generic layout. Endianness conversion is a bounded decode operation.
